Sources : études citées, présentées dans le complément « Les élasticités estimées sur données françaises » ; chaque chiffre, avec sa page et son tableau, figure dans donnees/elasticites\_france.csv.\par Lecture : chaque panneau ne rapproche que des études qui mesurent le même revenu, au même taux net, par la même méthode : variation d'une année sur l'autre, instrumentée par l'effet mécanique des réformes. Lehmann et al. et Sicsic exploitent les enquêtes Revenus fiscaux de l'Insee. Le rapport et l'article portent sur les foyers dont le revenu fiscal de référence dépasse 30 000 \euro{} ; l'article écarte en plus les variations extrêmes du revenu du capital.
